\section{Wall routing, guard barriers, and realization}\label{sec:walls}\label{wl:section}

Fix $2<p<\infty$.  We use the quantizer and carrier of
Proposition~\ref{qt:quantizer} and Lemma~\ref{qt:central-carrier}, the exterior
of Proposition~\ref{qt:exterior}, and the notation
\[
 M_y\subset\Lambda,\qquad M_x=h_*^{-1}(M_y),\qquad
 F_*=\overline K h_*,\qquad q_*=TF_*,\qquad q_0=T\widehat F.
\]
In particular $h_*$ is fixed, $q_*=Th_*$ near $\partial M_x$, and the
prescribed map on every full block is fixed.  Both exteriors are
irreducible and have vanishing absolute second homology.  The two
standard wall families are the annuli $A(D,t)$ and meridian disks
$B(e,s)$ of Proposition~\ref{qt:exterior}.  Their radial and perimeter
coordinates are
\[
 t=L_D\log r,\qquad
 q=d_D+r\bigl(\xi_e(s)-d_D\bigr),\qquad r\geq r_{u,D}>0.
\]
All families and coordinate decompositions below are locally finite.
A compact wall is understood as a proper wall in the manifold with
boundary $M_x$ or $M_y$.

This section separates the topological conclusion from its quantitative
inputs.  The pre-stack geometry is supplied in
Proposition~\ref{ms:prestacks}; the capacities used to choose guards are
fixed in Lemma~\ref{ha:capacities}, and
Proposition~\ref{ha:initial-walls} constructs the initial walls with
those bounds.  In particular, a large
Lipschitz constant for a subsequently chosen extension is never used as
an input to the guard construction.

\subsection{Sampling transverse labels}

Fix $0<\zeta<1$.  Let $j$ index a sample in one radial or perimeter
feature.  Write
$[0,w_j]$ for its root-parameter interval and
\[
 p_j:[0,w_j]\longrightarrow I_j
\]
for its increasing affine identification with a label interval about
the sample.  Within one feature the intervals $I_j$ are ordered and
disjoint, and $|I_j|\asymp c_*w_j$, where $c_*>0$ is fixed and small.
The number $w_j$ measures transverse parameter length, not physical
collar thickness.  Routing will retain a finite union
$E_j\subset(0,w_j)$ of closed parameter intervals with
\begin{equation}\label{wl:occupied-length}
 |E_j|\geq(1-\zeta)w_j.
\end{equation}
The same affine $p_j$ is used on every component of $E_j$.
We first construct stairs for arbitrary such sets.  Before the retained
intervals are chosen, we use the whole-interval case $E_j=[0,w_j]$
to test labels.  After realization, the stairs will be constant off
$p_j(E_j)$, so their nonzero derivatives occur only where a wall
parameter has been prescribed.

\begin{lemma}[Stairs and sampling]\label{wl:stairs}
On a feature interval of length $\ell$, one can choose $N$ equal
weights $w_j=w$ with $Nw/\ell=1+O(\zeta)$, and construct an
endpoint-fixing nondecreasing stair $S$, constant outside the selected
$p_j(E_j)$, such that its total increase over each $I_j$ is $\ell/N$.
Its slope with respect to the occupied root parameter is at most
$1+O(\zeta)$, and its Lipschitz constant in the target label is
$O(c_*^{-1})$.  The stair can be made arbitrarily close to the identity
on each compact feature.

The choices allow finitely many forbidden intervals of arbitrarily
small prescribed total length.  For a fixed finite list of integrable
edge-count tests, the weighted sample counts have expected density at
most $1+O(\zeta+c_*)$ against label measure.  They satisfy the
corresponding upper bounds with arbitrarily small additive errors and
probability tending to one as the slots are refined.
\end{lemma}

\begin{proof}
Remove the forbidden intervals and divide each remaining component
into slots.  Choose the numbers of slots approximately proportional to
component lengths.  Taking $w$ sufficiently small makes both rounding
errors and the difference of each slot length from $w$ as small as
required.  This also works for finitely many features using one common
$w$; no exact divisibility is required.  Sample in each slot outside a
small endpoint margin wide enough to contain $I_j$.  The total omitted
relative length and the rounding loss can be included in $\zeta$.

Give $S$ increase $\ell/N$ at constant speed over the occupied
root length $|E_j|$, and no increase on the complementary pieces.
Its speed is
\[
 \frac{\ell}{N|E_j|}\leq
 \frac{\ell}{Nw(1-\zeta)}=1+O(\zeta).
\]
Since $p_j'=|I_j|/w\asymp c_*$, the target-label slope is
$O(c_*^{-1})$.  Start the first flat portion at the left endpoint and
use all $N$ increments, so the final value is the right endpoint.
The ordered slots show that the displacement is bounded by their
largest length, the rounding error, and the total forbidden length.
These quantities can all be prescribed small.  The same conclusion
holds when $E_j=[0,w_j]$, a version used before selecting stable
occupied intervals.

For a nonnegative integrable count function $n$, sampling uniformly in
the allowed middle of a slot $J_j$ gives
\[
 \mathbb E\bigl[w_j n(\tau_j)\bigr]
 =\frac{w_j}{|J_j^{\rm allow}|}
       \int_{J_j^{\rm allow}}n(t)\,dt.
\]
The density is at most $1+O(\zeta+c_*)$.  For bounded $n$, independence
and $\max_jw_j\to0$ make the variance of the weighted sum tend to
zero.  For general $n\in L^1$, truncate it, apply this argument to the
bounded part, and use the first moment of the integrable tail.  A
finite list of tests can be imposed simultaneously.  On an exhaustion,
feature weights and additive errors can be chosen locally, with
summable failure allowances.
\end{proof}

\subsection{The two-colour pre-stack input}
\label{wl:prestack-subsection}

Before routing, a pre-stack is a product neighborhood of a sampled source
wall, equipped with its transverse scalar parameter.  Opposite colours may
still meet in circles as well as in the prescribed spanning arcs.  The
following definition records the simultaneous product structure and the
pointwise slope control used in routing and in the guard estimates.

\begin{definition}[Pre-stack contract]\label{wl:prestack}
Let $M_x,M_y$ be the source and target exteriors and let
$h_*:M_x\to M_y$ be the fixed reference homeomorphism.  A pre-stack
system consists of a locally finite, countable, two-colour family of
compact collars
\[
 \mathcal E_j:\Sigma_j\times[0,w_j]\longrightarrow C_j\subset M_x,
 \qquad u_j(\mathcal E_j(z,\lambda))=\lambda,\qquad w_j>0,
\]
whose central sampled wall is the level $u_j=w_j/2$, subject to the
following conditions.
\begin{enumerate}
\item Each $\mathcal E_j$ is a bi-Lipschitz homeomorphism onto its image.
For colour $A$, $\Sigma_j$ is an annulus whose core is essential in
$M_x$; for colour $B$, $\Sigma_j$ is a disk.
Each level is properly embedded, and
$C_j\cap\partial M_x=\mathcal E_j(\partial\Sigma_j\times[0,w_j])$.
The collars of a single colour are disjoint, including their boundaries.
The family of closed collars is locally finite.

\item There is a specified affine increasing correspondence from
$[0,w_j]$ to a true target label interval $I_j$.
On $\partial M_x$ every level has exactly the boundary trace of its
specified standard target wall, pulled back by $h_*$.
The increasing-$u_j$ coorientation agrees with that target label
coorientation at every boundary component; for an annulus this includes
both ends.  The target incidence says that an incident opposite-colour
pair has two boundary intersection points, one on each annulus end,
and that a nonincident pair has none.

\item For each opposite-colour pair $j,k$, the intersection
$C_j\cap C_k$ is a finite disjoint union of bi-Lipschitz product boxes
\[
 Q\times[0,w_j]\times[0,w_k],
 \qquad u_j=\lambda,\quad u_k=\mu,
\]
where $Q$ is a circle or a compact interval.  The interval endpoints
are exactly the parts of that box on $\partial M_x$.
These are full parameter rectangles: no box terminates at an interior
value of either parameter.  The level sheets meet transversely in
these product coordinates.  There are no further intersection pieces.
In particular, for an incident pair every level pair has the specified
single spanning arc, together with possible circles; a nonincident pair
has only possible circles.

\item In each collar the subdivision by all its opposite-colour boxes
is jointly bi-Lipschitz trivial over $u_j$, preserving every concurrent
opposite parameter.  Precisely, for some reference level there is a
bi-Lipschitz product parametrization that carries every intersection
box, every complementary piece, and every boundary piece at the
reference level to its counterpart at each level; on an opposite box
it preserves $u_k$.  Its restriction to a boundary piece remains in
$\partial M_x$.  No uniform quantitative bound for this
trivialization is required.

\item The ambient PL structure and the physical metric are compatible
in the following explicit sense.  Every locally finite arrangement
of pieces obtained by finitely many constant or affine parameter cuts
on each compact set admits common ambient coordinate changes that are
locally bi-Lipschitz in the physical metric and make the pieces PL
simultaneously.  These changes preserve the manifold boundary and
the specified incidences.  Separate unrelated PL charts for individual
pieces do not suffice for this requirement.

\item Every $u_j$ is Lipschitz on its closed collar in the physical
metric.  Any subsequent estimate using a density $G$ assumes, as an
additional quantitative input, a nonnegative locally measurable $G$
dominating the point-to-base upper local slopes
\[
 \operatorname{lip}_{C_j}u_j(x)
   =\limsup_{\substack{y\to x\\y\in C_j,\ y\ne x}}
         \frac{|u_j(y)-u_j(x)|}{|y-x|}.
\]
Use the adjacent maxima at internal interfaces and include both
parameter bounds at an essential intersection box.  On exceptional
collar-boundary or mesh-interface strata one may enlarge $G$ to the
actual local Lipschitz bounds of the fixed compact collars.  These
strata, in the common working coordinates or their bi-Lipschitz
images, have zero volume; the enlargement therefore changes no
volume capacity.  This convention does not require a two-point
Lipschitz bound across gaps of a nonconvex union of pieces.  The
contract alone asserts neither an integral estimate for $G$ nor
independence of that estimate from later choices of guards.
\end{enumerate}
\end{definition}

The full parameter rectangles and the simultaneous trivializations
let the pieces cut out by one colour be tracked as the other root
parameter varies, while preserving the concurrent opposite parameter.
The common ambient coordinates
serve a different purpose: they permit the eventual endpoint walls to
be cut and filled in one PL structure.  Constructions of these data,
including the sufficient test in Lemma~\ref{wl:pattern}, are given in
Sections~\ref{sec:mesh} and~\ref{sec:calibrated-islands}.

\subsection{Diagonal routing and conservation of periods}
\label{wl:routing-subsection}

The switching construction is related to regular exchange in
normal-surface theory~\cite{Haken1961}.  Here the additional requirement
is to preserve transverse parameter length and exact boundary labels,
including along a possibly infinite rooted assembly.
We work in the full source exterior $M_x$, including any boundary-collar
continuations used in preparing the input walls.  The constructions in
this subsection
assume the pre-stack contract of Definition~\ref{wl:prestack}.  They are
qualitative until explicit parameter-slope bounds are imposed.  We also use
the previously established irreducibility and vanishing
$H_2(M_x;\mathbb Z)=0$ of the exterior.

Fix one level in each of two intersecting opposite-colour collars.  A circle
of intersection bounds a disk on the disk wall.  It cannot be essential on
the annular wall, since an annular core is noncontractible in $M_x$.
Consequently every intersection circle is trivial on both walls.  A pair
with no prescribed boundary incidence has no intersection arc.  An incident
pair has exactly two boundary intersection points, one on each end of the
annulus, and hence exactly one proper intersection arc.  That arc spans the
annulus.  Any additional intersections are circles.

In a circle box write its coordinates as
\[
                (\theta,u,v)\in S^1\times[0,a]\times[0,b],
                \qquad a=w_j,\quad b=w_k.
\]
Reverse $u$ or $v$ in this box if necessary.  We arrange that the rim $v=0$
of the strip on a $j$-layer faces the exterior of its bounded inner disk,
and that the rim $u=0$ on a $k$-layer has the same property.  Call these two
sides OUT and the sides $v=b$ and $u=a$ IN.  The pattern trivializations
ensure that these choices hold throughout the box.  Introduce
\begin{equation}
\label{wl:diagonal-coordinate}
                              d=a-u+v.
\end{equation}
For $d\in(0,a+b)\setminus\{a,b\}$ the segment of this level in the parameter
rectangle has one endpoint on OUT and the other on IN.  Its product with
$S^1$ is an annulus.  On OUT the two sides give the $d$-intervals $[0,a]$
and $[a,a+b]$, while on IN they give $[0,b]$ and $[b,a+b]$.
Thus these annuli match the combined OUT parameter length to the combined
IN parameter length bijectively, by partial affine maps of slope $1$ or
$-1$.  The corner values $0,a,b,a+b$ are excluded.  No diagonal replacement
is made in an arc box: both original parameters are retained there.
Figure~\ref{wl:routing-figure} shows this parameter matching.

\begin{figure}[htbp]
\centering
\begin{tikzpicture}[x=1.1cm,y=1.1cm,>=Stealth]
 \fill[blue!4] (0,0) rectangle (4,3);
 \draw[thick] (0,0) rectangle (4,3);
 \draw[blue!65!black,thick,->] (3,0)--(4,1);
 \draw[blue!65!black,thick,->] (1.6,0)--(4,2.4);
 \draw[blue!65!black,thick,->] (0,.8)--(2.2,3);
 \draw[blue!65!black,thick,->] (0,2)--(1,3);
 \node[below] at (2,0) {OUT: $v=0$};
 \node[left,align=right] at (0,1.5) {OUT\\$u=0$};
 \node[above] at (2,3) {IN: $v=b$};
 \node[right,align=left] at (4,1.5) {IN\\$u=a$};
 \draw[->] (0,-.85)--(4,-.85) node[right] {$u$};
 \node[below] at (0,-.85) {$0$};
 \node[below] at (4,-.85) {$a=w_j$};
 \draw[->] (-1.15,0)--(-1.15,3) node[above] {$v$};
 \node[left] at (-1.15,0) {$0$};
 \node[left] at (-1.15,3) {$b=w_k$};
 \node[fill=white,inner sep=3pt] at (2.05,1.45) {$d=a-u+v$};
\end{tikzpicture}
\caption{The circle-box parameter rectangle.  Each regular $d$-segment
joins one OUT side to one IN side, with partial parameter transport of
absolute slope one.  Its product with the circle is the replacement
annulus.  Corner values are discarded; essential arc boxes retain both
parameters and are not routed.}
\label{wl:routing-figure}
\end{figure}


\begin{proposition}[Routing and good parameters]
\label{wl:routing}
For every pre-stack root $j$, outside a null subset of $(0,w_j)$, the diagonal
rule constructs a connected, oriented, properly embedded surface
$L_{j,\lambda}$ whose boundary, with coorientation, is the prescribed boundary
of the original $j$-level at parameter $\lambda$.  Different constructed
surfaces are disjoint except for the prescribed essential arc intersections
between opposite-colour roots.  Each such surface is obtained by attaching
punctured disks and annuli along a rooted tree; it may have ends at infinity.
The following stronger local finiteness holds for each constructed surface:
for every compact $K\subset M_x$, only finitely many of its reached states
have an original parent collar meeting $K$.
\end{proposition}

\begin{proof}
Cut every original level along the open strips belonging to circle boxes.
The circles are trivial and disjoint on that level.  Their bounded disks
are therefore nested or disjoint.  There is one root piece containing the
actual boundary of the wall, and, for an annulus, both actual boundary
components belong to this same piece.  The root is an annulus or a disk with
finitely many open disks removed.  Every other piece is a punctured disk.
Each nonroot piece has exactly one IN rim, namely its rim toward its parent
in the nesting, and any remaining rims are OUT.  The root has no IN rim.
All essential arc boxes lie on root pieces: a spanning arc or a proper arc
with endpoints on the actual boundary cannot enter a bounded disk through
one of its disjoint circle boundaries.  There are finitely many pieces in
any one collar, and they are tracked over its whole parameter interval by
the pattern trivialization.

A state is a pair $(P,t)$ consisting of such a tracked piece and its local
parameter.  Starting at the root state $(R_j,\lambda)$, follow every OUT rim
through its diagonal annulus to the matched IN rim, and continue at that
state.  Every nonroot state has exactly one incoming predecessor wherever
the matching is defined: its unique IN rim has precisely one OUT partner.
Root states have none.  This proves that no state reached from a root can
repeat.  Indeed a repetition along an ancestral path would give a cycle of
predecessors, which could never terminate at the root.  If two distinct
ancestral paths reached the same state, following its unique predecessors
backwards would force their paths, root indices, and root parameters to
coincide.  This also rules out duplication between different root-parameter
starts.  The states reached from one start consequently form a tree, with
finite branching at each state.  Cycles or infinite backward orbits that
are not reachable from any root play no role.

There are countably many tracked pieces and countably many finite transition
itineraries.  On its domain each itinerary maps the starting parameter to
a local parameter by $t=\lambda+c$ or $t=-\lambda+c$.  A corner or breakpoint
condition therefore excludes at most one starting value for that itinerary.
Discard the union of these countable sets.  All transitions then occur along
nondegenerate segments and are local partial isometries of parameter
intervals.

Here is the estimate that controls possibly infinite trees.  For a tracked
piece $P$ belonging to collar $C_k$, let $N_P(j,\lambda)$ count the states
of the tree started at $(j,\lambda)$ that use $P$.  The image parameter sets
in $(0,w_k)$ of all valid itineraries from all roots to $P$ are pairwise
disjoint.  Otherwise one state would have two backward ancestries.  Each
itinerary preserves one-dimensional Lebesgue measure.  Decomposing its
valid domain into its measurable partial intervals and using countable
additivity gives
\begin{equation}
\label{wl:visit-bound}
             \sum_j\int_0^{w_j} N_P(j,\lambda)\,d\lambda\leq w_k.
\end{equation}
The valid domains are measurable: finite itineraries impose affine interval
conditions, and the excluded parameter set is countable.

For a compact $K$, let $\mathcal P(K)$ be all tracked pieces whose original
parent collars meet $K$.  This set is finite, by local finiteness of the
compact collars and the finite number of pieces in each collar.  Summing
\eqref{wl:visit-bound} yields
\[
 \sum_j\int_0^{w_j}
      \sum_{P\in\mathcal P(K)}N_P(j,\lambda)\,d\lambda
       \leq\sum_{P\in\mathcal P(K)}w_{k(P)}<\infty.
\]
It follows that for almost every parameter of every root, only finitely
many reached states have parent collars meeting $K$.  Intersect these
full-measure sets over a countable compact exhaustion of $M_x$.  Call the
remaining parameters good, also requiring avoidance of all corner values.

For a good parameter, glue all reached pieces to the associated diagonal
annuli.  The gluing is locally finite in the ambient manifold: a band meeting
a compact set lies in a circle box, hence in the original collar of its
parent state, and those states are finite in number by goodness.  Each
piece or band is compact and includes its attaching rims.  Thus the assembly
is locally closed and properly embedded.  Each rim has exactly two local
half-surface pieces, giving an ordinary surface chart there.  The pre-stack
charts make the same assertion at corners and at the actual manifold
boundary.  Connectedness follows from the tree construction.  Only the root
piece meets the actual boundary, so the boundary of the assembly is exactly
the original root boundary.

Orient the root according to its prescribed coorientation in the oriented
three-manifold.  At each tree edge orient the attached band and then its
child piece so that the common boundary orientations cancel.  There is no
orientation consistency obstruction because no state is attached by two
ancestral paths.  All child pieces are planar and all bands are annuli.
A finite tree assembly, capped at its frontier rims by disks, is therefore
a disk or annulus of the original root type: removing an inner disk and
attaching a punctured disk followed by disks at its remaining holes simply
replaces that inner disk by a disk.

Finally, two assemblies cannot meet in one original piece at the same
parameter or in one diagonal annulus, by uniqueness of backwards ancestry
and the bijective diagonal matching.  Their only possible remaining
intersections are in arc boxes.  Those boxes belong to root pieces on both
sides and are exactly the retained essential intersections.  A single
assembly contains only one root, so these intersections cause no
self-intersection.
\end{proof}

\begin{remark}
\label{wl:routing-no-width-arithmetic}
The argument uses ordinary Lebesgue length in each original parameter
interval and partial isometries between those intervals.  The widths $w_j$
need not be equal or commensurable.  Goodness controls states whose
\emph{entire original collars} meet a compact set; controlling only the
routed fragments would not suffice for the cap argument below.
\end{remark}

The routed surface has the right boundary, but it need not have been
obtained by a compact isotopy.  To compare its crossings with those of
the standard wall, we truncate the routing tree and cap its frontier.
This reduces every compact path test to a compact homology calculation.

\begin{proposition}[Conservation of oriented periods]
\label{wl:period}
Let $L=L_{j,\lambda}$ be a good routed surface.  Let $W_{j,\lambda}\subset M_x$
be the pullback by $h_*$ of the standard cut at the exact affine label
corresponding to $\lambda\in(0,w_j)$.  Give both surfaces the coorientation of
increasing standard parameter.  For every compact path $\gamma$ with endpoints
on $\partial M_x\setminus\partial L$, their relative oriented intersection
numbers agree:
\begin{equation}
\label{wl:period-identity}
                  I(\gamma,L)=I(\gamma,W_{j,\lambda}).
\end{equation}
Here intersection numbers may be defined after general-position perturbation
relative to the endpoints.  If the original path has finitely many hits on
$L$, its signed local crossings give the same number and
\[
                       |I(\gamma,L)|\leq\#(\gamma\cap L),
\]
where hits are counted with their occurrences along the path.
\end{proposition}

\begin{proof}
Truncate the routing tree at a finite depth $n$.  At every frontier OUT rim,
cap with the original inside disk bounded by that rim in its original wall
at the corresponding local parameter.  Orient each cap to cancel the
frontier boundary.  The result is a finite compact oriented two-chain $A_n$.
The caps need not be disjoint from one another or from the truncated assembly;
only their boundary identities are used.  Every cap is contained in the
original parent collar of its frontier state.

Let $K$ be a compact neighborhood of the path, or of any prescribed compact
homotopy of paths relative to their endpoints.  Goodness says that only
finitely many reached states have original collars meeting $K$.  Their
depths have a finite maximum.  If $n$ exceeds that maximum, every omitted
piece and every new cap lies outside $K$.  Taking a slightly larger compact
neighborhood first shows that $A_n$ agrees with $L$ in a neighborhood of the
path or homotopy.  This is the reason for using whole parent collars in the
definition of goodness.

The actual boundary of $A_n$ is the root boundary.  The pre-stack boundary
contract, including agreement of coorientations at both annular ends, gives
\[
                     \partial A_n=\partial W_{j,\lambda}
\]
as oriented one-chains.  A common subdivision and boundary parametrization
makes these chains literally identical.  Hence
$Z_n=A_n-W_{j,\lambda}$ is an absolute compact two-cycle.  Since
$H_2(M_x;\mathbb Z)=0$, it bounds a finite compact three-chain $B_n$.

We spell out why the endpoints on the manifold boundary cause no extra
intersection term.  First take the path in the interior except at its
endpoints, using a boundary collar and keeping the endpoints fixed.  Neither
endpoint lies on the support of $Z_n$ at the boundary, since the two surface
boundaries agree and the endpoints avoid that common boundary.  Push the
cycle and its bounding chain into the interior by a sufficiently small
collar push.  During this push the cycle avoids neighborhoods of the two
endpoints, so its intersection with the path is unchanged.  The pushed
three-chain is disjoint from the path endpoints.  The boundary formula for
oriented intersections, after general position, therefore gives
\[
                  I(\gamma,Z_n)
                    =I(\gamma,\partial B_n)=0.
\]
The same formula applies to the pushed chains; the notation suppresses that
harmless push.  Since $A_n=L$ near the path, this proves
\eqref{wl:period-identity}.  For a homotopy relative to the endpoints, choose
one $n$ for its entire compact image.  The same finite-chain argument proves
homotopy invariance.  Equivalently, one can push that homotopy into the
interior away from the fixed endpoint neighborhoods before applying the
intersection formula.

If the unperturbed path has finitely many hits, choose disjoint small
parameter intervals about those hits, each mapping into a two-sided surface
chart.  Before and after its isolated hit the path lies on a definite local
side.  Its local algebraic crossing is $1$, $-1$, or $0$, according to these
two sides and the coorientation.  A sufficiently small general-position
perturbation has exactly that total algebraic crossing in the chart.  The
same reasoning works in a proper half-space chart at a boundary hit; push
its path segment slightly inward and keep its ends fixed.  Outside these
intervals there are positive gaps on compact remainders, so no new crossings
are needed.  Summing proves the asserted agreement and the bound by the
number of hits.  No isotopy relative to any later guard system has been
used.
\end{proof}


\subsection{Graph potentials and guards}

On $M_y$ put $\beta=\pi T\in\Gamma$, with subedge lengths inherited
from Euclidean arclength.  Define $\alpha$ in a metric star having one
branch for each subcell $D$, with branch coordinate $b=-t\geq0$ and
common node $b=0$.  Thus all $r=1$ points have the same $\alpha$.
The coordinates are continuous and locally Lipschitz in the product
charts.  If $q=T(y)$, then
\begin{equation}\label{wl:graph-physical}
 |q-\beta|\leq Cb.
\end{equation}
At a fixed compatible $\beta$ on one branch, the physical point varies
with $b$ at speed at most $C$; at fixed $b$ in one subcell it varies by
at most the Euclidean displacement of $\beta$.  These follow directly
from the polar formula and $L_D\asymp\max|\xi_e-d_D|$.

For a layer rooted at $A(D,t_L)$ write $\alpha_L$ for its standard
star label.  For a layer rooted at $B(e,s_L)$ write $\beta_L$ for its
standard perimeter label.  The following potentials, each
$1$-Lipschitz in the appropriate graph metric, detect distance from
the entire physical support of the root layer:
\begin{equation}\label{wl:potentials}
\begin{array}{c|c|c}
 \hbox{root}&\hbox{star potential}&\hbox{perimeter potential}\\ \hline
 A(D,t_L)&\dist(\alpha,\alpha_L)&\dist_{\R^3}(\beta,\partial D)\\
 B(e,s_L)&\dist(\alpha,\mathcal S_e)&|\beta-\beta_L|.
\end{array}
\end{equation}
Here $\mathcal S_e$ is the subtree consisting of the common node and
all full branches corresponding to subcells incident to $e$.

To verify detection, suppose both displayed distances are small.  In
the $A$ case, if the guard is on the same branch, compare at a nearby
perimeter point of $\partial D$ and then at the same $b$.  On another
branch, small star distance implies that both branch heights are
small, so \eqref{wl:graph-physical} applies.  In the $B$ case, on a
compatible branch compare using the same $b$ and $\beta_L$; on an
incompatible branch small distance to $\mathcal S_e$ makes the guard
height small.  Thus the physical distance to the standard support is
at most a fixed multiple of the sum of the two potential values.
Both potentials vanish on that support.  Moreover, when the potential
tests the colour opposite to the root, it vanishes identically along
every feature having an essential intersection with the root.

\begin{definition}[Guards and prior capacities]\label{wl:guards}
Choose a locally finite compact cover of $M_x$ by smaller buffered
coordinate regions $K_a$, with larger regions $U_a$ carrying fixed
bi-Lipschitz ball or half-ball charts.  All smaller regions have
positive margins in the larger charts.  Fix finite numbers $C_a$ as
upper capacities for
\begin{equation}\label{wl:prior-capacity}
 \int_{U_a}G^p\leq C_a.
\end{equation}
These numbers, the coordinate bounds and the prescribed continuous
label accuracies are fixed before guards, sample weights and the
final resolution are chosen.

A guard system is obtained from a sufficiently fine homeomorphic
approximation $h_g$ to $F_*$, agreeing with $h_*$ on the exact boundary
data and respecting the exteriors.  Choose a locally finite source
net of sufficiently small local spacing.  For every net point $g$,
choose two paths in the constant $Th_g$ fibre from $g$ to boundary
anchors whose $h_*$-images are extremal points of the target fibre
on distinct full blocks, and let $\kappa_g$ be their union.  Thus
$\kappa_g$ is compact and connected, and each of its points has a
path within $\kappa_g$ to a boundary anchor.  Write
$q_g=Th_g(g)$ and require $Th_g$ to be sufficiently close to $q_*$.
The pre-stacks must be disjoint from all $\kappa_g$.
\end{definition}

The guards have a locally uniform positive reach, independent of the
net spacing and the approximating $h_g$.  Indeed a target fibre in
the exterior has paths to at least two extremal endpoints on distinct
full blocks, also over true edges and vertices.  Those endpoints have
positive mutual separation on compact label tests.  Their source
positions use the fixed boundary map $h_*^{-1}$, so a guard through
any one point has uniformly positive diameter.  Properness of
$Th_g$ makes the family of guards locally finite: if a guard meets a
fixed compact set, its label belongs to a compact label set, and its
net point belongs to the compact inverse image of that set.

For later use, the order in Definition~\ref{wl:guards} can be
implemented without a relative isotopy fixing guards.  Reserve tiny
forbidden parameter neighborhoods of the guard labels in each
feature before sampling.  Their total length can be arbitrarily
small.  Properness bounds the number of relevant guards for each
compact feature using the prescribed net; hence the accuracy of
$h_g$ can be chosen before the actual gaps are placed.  Carrier
levels outside slightly enlarged gaps have positive buffers from the
guards.  Make the mesh and boundary motions smaller than these
buffers.  The preparation in Proposition~\ref{wl:normalization} ends with
wall pieces lying only in tetrahedra incident to original possible
edge hits.
Controlled placement, thickening and cell-sized snapping then stay
within the buffers.  This controls the final wall supports; no
normalization isotopy relative to the guards is required.
This argument is useful only when
\eqref{wl:prior-capacity} remains valid at all such smaller
resolutions; that is the independent capacity assertion of
Lemma~\ref{ha:capacities}, with the geometry of
Proposition~\ref{ms:prestacks}.  Proposition~\ref{ha:initial-walls}
carries out the preparation with these choices.

\begin{lemma}[Crossing lower bound]\label{wl:crossing-bound}
Assume the pre-stack hypotheses of Definition~\ref{wl:prestack} and
its slope majorant $G$, and let the pre-stacks avoid the guards.
Suppose a guard label has physical distance at least $\delta>0$ from
the standard support of a good routed layer $L$.  On a fixed compact
separation test, sufficiently accurate feature stairs give a
constant $c_\delta>0$ such that every compact rectifiable path $\sigma$
joining $L$ to that guard satisfies
\begin{equation}\label{wl:path-cost}
 \int_\sigma G\,d\mathcal H^1\geq c_\delta.
\end{equation}
The constant depends on the fixed geometry and separation test, not
on the sample weights or the number of switches in an itinerary.
\end{lemma}

\begin{proof}
We may assume $\int_\sigma G\,d\mathcal H^1<\infty$; otherwise
the conclusion is immediate.  The coarea bound below will make the
weighted crossing sums absolutely integrable.

Use the whole-interval stairs of Lemma~\ref{wl:stairs}, before
restricting to stable occupied intervals.  They act on the star and
on $\Gamma$, fixing vertices and each edge setwise; on the star use
$b=-t$ with its orientation reversed.  Choose one of the potentials
$U$ in \eqref{wl:potentials} giving a separation comparable to
$\delta$.  For each root $j$ of the tested colour set
\begin{equation}\label{wl:crossing-coefficient}
 c_j(\lambda)=\frac{d}{d\lambda}
                 U\bigl(S(p_j(\lambda))\bigr).
\end{equation}
These coefficients are uniformly bounded.  In the opposite-colour
case they vanish almost everywhere on all essentially incident
features.  The potential at the boundary anchor of $L$ is as close
to zero as required, whereas the potential at the guard anchor stays
quantitatively positive.

Extend $\sigma$ along $L$ and along the guard to a path $\gamma$
from a boundary anchor $a_L$ of $L$ to a boundary anchor $a_g$ of
the guard.  For almost every tested root parameter with nonzero
coefficient, the added portions avoid the corresponding routed
layer: the guard avoids all pre-stacks, distinct routed layers are
disjoint except at essential root boxes, and coefficients on those
excluded essential features vanish.  The layer containing the
chosen starting point itself is a single null parameter value.
Let $\ell$ denote the tested graph coordinate, either $\alpha$ or
$\beta$.  The weighted crossing identity is
\begin{equation}\label{wl:weighted-period}
\begin{aligned}
 \sum_j\int_0^{w_j}c_j(\lambda)I(\gamma,L_{j,\lambda})\,d\lambda
 &=\sum_j\int_0^{w_j}c_j(\lambda)I(\gamma,W_{j,\lambda})\,d\lambda\\
 &=U\bigl(S(\ell(h_*(a_g)))\bigr)
       -U\bigl(S(\ell(h_*(a_L)))\bigr).
\end{aligned}
\end{equation}
Here the sum is over roots of the tested colour.  The first equality
is Proposition~\ref{wl:period}.  For the second, choose a fixed
piecewise regular representative of $\gamma$, homotopic relative
to its endpoints.  It meets only finitely many standard features;
its graph-coordinate
paths are piecewise Lipschitz.  Oriented one-dimensional level
counting and the Fundamental Theorem of Calculus on every graph
edge give the displayed difference of $U\circ S$ at the two anchors.
There is no error accumulated along a cycle or a long path.

On the other hand, one-dimensional coarea bounds the integrated
number of hits on $\sigma$.  Within an unchanged pre-stack piece it
uses $u_j$; within a transition rectangle it uses
$d=w_j-u+v$.  The upper local slope of $d$ is at most the sum of the
two original slopes.  Partial transition itineraries have derivative
$+1$ or $-1$, and Proposition~\ref{wl:routing} gives a unique reachable
ancestry for every piece-value and every diagonal-segment value.
Thus coarea pays for all reached root layers there without an
itinerary multiplicity.  Same-colour disjointness bounds overlap;
at an essential box count both parameters.  Apply coarea on the
closed subsets of path portions and partition the affine itinerary
rules countably if necessary.  Interface slopes use the convention in
Definition~\ref{wl:prestack}.  More explicitly, the fixed parameter
function is Lipschitz on each closed collar.  On the closed subset
of times when a rectifiable path lies in that collar, its scalar
trace has a Lipschitz extension; at almost every density time its
absolute derivative is bounded by the point-to-base upper slope
times the path speed.  The one-dimensional coarea inequality on
that subset therefore uses exactly this $G$.  Hence
\[
 \sum_j\int_0^{w_j}
       \#\bigl(\sigma\cap L_{j,\lambda}\bigr)\,d\lambda
       \leq C\int_\sigma G\,d\mathcal H^1.
\]
The right side is finite by assumption.  Hence almost every count
is finite and the weighted signed sum in \eqref{wl:weighted-period} is absolutely
integrable.  The appended portions contribute zero for almost every
parameter with nonzero coefficient, so the displayed coarea bound
controls its absolute value.  The bounded coefficients and the
positive potential difference prove \eqref{wl:path-cost}.

There is no infinite list of conflicting stair accuracies hidden in
this argument.  The perimeter projection is proper.  Near a compact
guard perimeter-label test only finitely many subcells and subedges
occur, and they use finite lists of accuracy requirements.  For a remote $A$
feature use distance to its whole $\partial D$; for a remote $B$
feature use distance to its whole closed subedge.  These supports are
invariant under their stairs, and only stair accuracy near the guard
label is needed.  The proper map $\pi q_*$ sends a locally finite
compact source cover to locally finite compact perimeter-label sets.
Choose locally finite, relatively compact enlargements of these sets
and accuracies smaller than their margins.  This supplies consistent
local prescriptions on all features.
\end{proof}

\begin{lemma}[Spherical path estimate]\label{wl:sphere}
Let $p>2$, and let $\Sigma_R$ be a sphere of radius $R$, either whole
or intersected with a closed half-space containing its center.  There
is a constant $C_p$ such that, for any two points $z_1,z_2\in\Sigma_R$
and nonnegative measurable $g$ with finite surface $L^p$ norm, some
rectifiable path in $\Sigma_R$ joining them has cost at most
\begin{equation}\label{wl:sphere-upper}
 C_p R^{1-2/p}\left(\int_{\Sigma_R}g^p\,d\mathcal H^2\right)^{1/p}.
\end{equation}
Fixed bi-Lipschitz chart changes alter only the constant.
\end{lemma}

\begin{proof}
Scale first to $R=1$.  Choose a fixed cap strictly above the equator
parallel to the cutting plane.  From an endpoint on or above the
equator, shortest arcs to cap points stay in the upper hemisphere.
For an endpoint below the equator, restrict cap points to those with
nonnegative horizontal scalar product with its radius direction.
This retains a fixed fraction of the cap.  The initial vertical
velocity of the shortest arc is then nonnegative.  While its height
is negative, the geodesic equation gives nonnegative second
vertical derivative; once it becomes positive it cannot return
negative before its positive endpoint on an arc shorter than $\pi$.
The arc therefore stays above the height of its starting point and
inside the prescribed half-space.

Average the cost of these arcs over the permitted cap fraction.
Spherical polar integration from the endpoint gives a kernel bounded
by $C/\sin\theta$ against surface measure.  Its conjugate power is
integrable precisely because $p'=p/(p-1)<2$; the endpoint and
antipodal singularities both satisfy
$\int_0^\epsilon\theta^{1-p'}\,d\theta<\infty$.
Hölder's inequality bounds the averaged cost by $C_p\|g\|_{L^p}$.
Do this at both endpoints.  Join the resulting two cap points by a
shortest arc in the cap's containing upper hemisphere; averaging
also this middle cost has the same bound.  Some choice of the two
cap points has total cost bounded by the sum of the averages.
Scaling surface measure and length restores the factor
$R^{1-2/p}$.  The same proof bounds the costs after a fixed chart
change, using its length and surface-measure constants.
\end{proof}

\begin{proposition}[Guard barrier]\label{wl:barrier}
Fix the data of Definition~\ref{wl:guards}, the capacities $C_a$ and
a positive continuous local layer-support accuracy.  There are
sufficiently fine net spacings and sufficiently small approximation
and stair errors, depending only on these prior data, with the
following property.  If a pre-stack system satisfies
\eqref{wl:prior-capacity} and avoids the resulting guards, then every
point of every good routed layer is within the prescribed local
accuracy, under $q_*$, of that layer's standard physical support.
\end{proposition}

\begin{proof}
It suffices to impose countably many locally finite compact tests
with positive accuracy constants.  Suppose one fails: at a point
$x$ of a good layer, its physical-support discrepancy exceeds
$\delta$.  The layer is connected to its actual boundary, where its
label is exact.  Uniform continuity of $q_*$ on a buffered chart
therefore gives a path stretch of the layer of diameter at least
$a>0$ on which the discrepancy exceeds $\delta/2$.  The number $a$
can be chosen from the fixed smaller-chart margins, the modulus of
continuity and $\delta$, independently of the layer and sampling.
This also applies near the manifold boundary, using half-ball
charts; a discrepant point cannot end its path in the exact boundary
without leaving this discrepant neighborhood.

For small $h$, pack at least $c a/h$ disjoint chart balls of radius
$2h$ centered along this stretch.  A maximal separated selection on a
connected set of diameter $a$ gives this elementary packing bound.
Place a guard net point within $h/2$ of each center.  Its label is
still separated by a fixed fraction of $\delta$ from the layer
support, by uniform continuity and the chosen approximation
accuracy.  The layer stretch and the guard each meet every chart
sphere of radius between $h$ and $2h$ about the center: each has a
point inside that radius and a connected continuation outside it.
For the guard this uses the prior uniform reach; for the layer it
uses the diameter of the selected stretch and smaller-chart
margins.

Every spherical path between such two hits has $G$-cost at least a
fixed $c_\delta>0$ by Lemma~\ref{wl:crossing-bound}.  Apply
Lemma~\ref{wl:sphere} to the whole or truncated sphere.  For almost
every radius $R\in[h,2h]$,
\[
 \int_{\Sigma_R}G^p\,d\mathcal H^2\geq cR^{2-p}.
\]
Integrating in radius gives at least $c h^{3-p}$ in each chart ball.
The disjoint balls therefore force
\begin{equation}\label{wl:barrier-divergence}
 \int_{U_a}G^p\geq c a h^{2-p}.
\end{equation}
Since $p>2$, choose $h$ using the previously fixed $C_a$ to contradict
\eqref{wl:prior-capacity}.  All constants used before choosing $h$
come from the fixed test, not the pre-stack resolution or weights.
Spheres whose centers are near the boundary are exactly the
half-space intersections in Lemma~\ref{wl:sphere}.  Paths are allowed
in the whole manifold chart, including through guards; no estimate
on a punctured chart is used.  A locally finite choice of spacings,
with positive continuous minorants for the requested errors, proves
all the compact tests simultaneously.
\end{proof}

\subsection{Stable occupied stacks and their realization}

Choose the layer-support errors of Proposition~\ref{wl:barrier} also
small enough toward infinity to have two consequences.  A layer
with a fixed compact standard support is confined to a compact
source set; and on any fixed compact source set only finitely many
root indices can occur.  To make these choices explicitly, take a
compact exhaustion of the locally finite label complex, enlarge
each member by a compact neighborhood, and use properness of $q_*$
to pull these neighborhoods back.  Require the local error outside
the corresponding inverse image to be less than its positive
separation from the smaller label compact.  Conversely, on a compact
source test choose the error small enough to keep any meeting
standard support in a fixed compact enlargement of its $q_*$ image.
Only finitely many compact features, and finitely many samples on
each feature, occur there.

\begin{lemma}[Stable intervals]\label{wl:stable}
Under these confinement choices, a good routed layer is a compact
disk or annulus of its root type.  For every root $j$ and every
$\zeta>0$ one can choose finitely many separated closed intervals
$E_j\subset(0,w_j)$ satisfying \eqref{wl:occupied-length}, such that
their routed layers form compact, locally finite, bi-Lipschitz
product stacks.  These products preserve all concurrent essential
parameters and the exact boundary data.  Different stacks are
disjoint except for the genuine essential-pair products.
\end{lemma}

\begin{proof}
A good tree visits only finitely many states whose entire initial
collars meet a given compact set.  Confinement therefore makes the
whole reached tree finite.  It has no frontier, and attaching its
finite tree of punctured disk pieces merely fills the root's inner
disk holes by disks.  Thus its completed surface has the original
disk or annulus topology and is compact.

All corner and breakpoint values have been discarded.  Finitely
many strict affine itinerary inequalities then remain valid on an
open interval about the starting parameter.  The same finite tree
and switching sides persist throughout this interval.  The good
parameters have full measure, so a finite collection of separated
closed intervals inside these stability neighborhoods captures at
least $(1-\zeta)w_j$ of the root interval.  Locally finite confinement
persists by continuity from the dense good parameters.

For product triviality use the pre-stack trivializations under the
affine parameter transports on each reached piece.  In a circle
box, parametrize each nonvanishing diagonal segment between its
fixed pair of sides by fractional affine position.  Its product
with the circle gives a band bundle.  Match the boundary-circle
parametrizations of this band to the adjacent piece
trivializations: relative to a reference parameter, lift the
orientation-preserving circle reparametrizations continuously to
monotone maps of $\R$ commuting with integral translation, and
interpolate their lifts across the band.  On a compact stability
interval both the maps and their inverses have finite Lipschitz
bounds, as do the segment fractions.  Gluing the finitely many
bundles gives a bi-Lipschitz product.  Essential boxes lie on the
root piece and their other parameter is preserved by the original
trivialization, so the modifications do not alter them.  The
compatible changed-coordinate PL hypothesis in
Definition~\ref{wl:prestack} applies to all endpoint arrangements.
\end{proof}

\begin{theorem}[Conditional wall realization]\label{wl:realization}
Fix the quantizer, carrier, exterior and boundary data above.  Fix
$0<\zeta<1$, $c_*>0$, the requested locally positive label and value
accuracies, and locally finite buffered charts with finite
capacities $C_a$.  Choose guards and feature-stair accuracies in the
order of Definition~\ref{wl:guards} and
Proposition~\ref{wl:barrier}.  Suppose the resulting sampled
pre-stacks satisfy Definition~\ref{wl:prestack}, avoid these guards,
and admit a slope majorant $G$ obeying the previously fixed
capacities \eqref{wl:prior-capacity}.  The sample intervals are
ordered, have the affine parametrizations $p_j$, and satisfy the
slot, gap and accuracy choices of Lemma~\ref{wl:stairs}.

Then there are finite unions of occupied closed intervals $E_j$
satisfying \eqref{wl:occupied-length} and a locally bi-Lipschitz
homeomorphism
\[
 h:M_x\longrightarrow M_y,\qquad h|_{\partial M_x}=h_*|_{\partial M_x},
\]
which sends every occupied routed layer at parameter $\lambda$ to
the standard layer at label $p_j(\lambda)$, with the same affine
parametrization on all occupied components.  It extends by the
prescribed $h_*$ on the removed cores and tubes.  Moreover
$q_h=Th$ is as finely close to $q_*$ on $M_x$ as prescribed.

The Theorem uses only the preassigned capacities for its choice of
guards.  Existence of a system meeting its hypotheses at that
resolution is the separate content of
Proposition~\ref{ms:prestacks}, Lemma~\ref{ha:capacities}, and
Proposition~\ref{ha:initial-walls}.
\end{theorem}

\begin{proof}
Lemma~\ref{wl:stable} supplies the occupied stacks.  We first construct
$h$ with their exact parameter correspondence and then prove the
asserted label accuracy.

Cut at both endpoints of every occupied $B$ interval, thus at
multiple gates on every subedge.  The corresponding target
components are balls by the graph-handlebody description in
Proposition~\ref{qt:exterior}; this includes the intermediate
chambers along a subedge.  A target ball boundary consists of a
punctured-sphere patch of $\partial M_y$ and its distinct gate disks.
At the source the corresponding fixed boundary patch, capped with
the new gates, is an embedded sphere.  Trim a small product collar
of the manifold boundary and carry this sphere into the trimmed
copy.  It bounds a ball by irreducibility.  The ball lies on the
inward side: the discarded collar escapes to the original manifold
boundary without meeting the sphere, so it cannot be contained in
the ball.  Pulling back from the trimmed copy gives the required
ball bounded by the original patch and gates.

No additional gate lies inside this ball, since its boundary lies
outside the prescribed patch and it is disjoint from the ball's
gates.  Along every gate, the two prescribed patch-balls occupy its
two local sides, as already seen near its rim.

The patch-ball interiors are disjoint.  Perform the following
comparison in the preceding trimmed construction.  For any two
patch-balls, choose a point just inside each of their distinct
boundary patches, lying outside the other ball.  If the capped
spheres share gates, their opposite local sides permit arbitrarily
small compatible collar pushes there that make the spheres disjoint
while preserving these two points.  If their interiors overlapped,
the pushes could also retain a point of that overlap.  The resulting
disjoint spheres would then bound nested balls, contradicting one of
the two distinguished points.  Hence the original interiors cannot
overlap.

There is no leftover component.  Indeed each patch-ball minus its
gates is a relatively
open and closed connected region off the gate system.  A compact
transverse path from an unassigned component to an outer boundary
patch would have finitely many gate crossings by local finiteness;
it can neither have a first entry through a gate, whose two sides
are already assigned, nor reach the endpoint without such an
entry.  Thus the source and target $B$-ball decompositions
correspond exactly.

On each gate, the endpoint $A$ incidences are disjoint proper arcs
with prescribed endpoint pairs.  They correspond relative to the
circle boundary by ordinary cutting of a disk along proper arcs.
Cutting an endpoint annulus along all its spanning gate arcs leaves
disks.  Their boundary slots on the two annulus ends identify the
pieces, even if the spanning arcs twist.  There are at least two
distinct cuts on every sampled subedge: each occupied interval has
two endpoints and each feature is sampled.  Thus successive arcs
leave actual disk pieces with slots on both ends; with more than
two arcs the slot not containing another endpoint uniquely
identifies adjacent pieces.  Along a gate interior there is one
adjacent piece on each local side, in the two distinct gate balls.

These $A$ disks lie in the matching $B$ balls, as their end segments
already do.  Their boundary loops correspond on the ball spheres.
A disjoint family of proper tame disks in a ball cuts it into balls,
indexed by the sphere patches cut along those loops.  Each such
patch reaches an open patch of the manifold boundary: on a gate,
every complementary region of disjoint proper arcs reaches an open
rim interval.  Consequently the endpoint chamber incidence agrees
on the source and target.  The collar structure makes membership
in the interior of an occupied interval constant on an endpoint
cell interior, and its boundary patch determines the corresponding
active or inactive side.

All of these disk and ball recognitions can be made in the common
changed polyhedral coordinates required by
Definition~\ref{wl:prestack}.  Their finite PL models are
bi-Lipschitz equivalent to ordinary disks and balls.  Tame
Schoenflies and the relative PL structure are the qualitative
topological inputs; see \cite[Theorem~5]{Brown1960} and
\cite[Theorem~2.2]{Hamilton1976}.  No uniform constant for these
models is claimed or required.

It remains to map the occupied interiors, not just their endpoint
walls.  A region where both colours are active is an arc times the
two occupied parameter intervals.  Its two boundary end patches
already carry the prescribed correspondence $h_*$.  Extend this as
a product, interpolating the interval coordinate along the arc
between its two increasing endpoint correspondences, and preserve
the two root parameters exactly.  In a $B$ stack, removing the
interiors of these opposite-colour strips from each disk fibre
leaves disk fibres.  In an $A$ stack the analogous removal of
spanning strips leaves disk fibres as well.  Their incidence was
identified above, and Lemma~\ref{wl:stable} gives products over the
remaining single active parameter that preserve all concurrent
box parameters.

The boundaries of each such disk fibre are already assigned, either
on a both-active product or on $\partial M_x$.  In disk-product
coordinates extend them by coning.  More explicitly, for a jointly
bi-Lipschitz family of circle maps $g_\lambda$, the map
\[
 (r\theta,\lambda)\longmapsto
             (r g_\lambda(\theta),\lambda)
\]
is jointly bi-Lipschitz, as is its inverse; the factor $r$ controls
the parameter variation at the disk center.  The boundary circle
maps have these bounds because their finitely many pieces and the
product charts do.  This defines the single-active portions
compatibly with the both-active ones.

These active products occupy the appropriate endpoint cells
completely: their fibres reach the true boundary, their side walls
are precisely the prescribed cuts, and the collar/product
structure makes them open in the relevant cell interiors and
closed up to their indicated boundary.  Every remaining open cell
is therefore a neither-active ball.  Its boundary map is already
assigned and extends by coning in bi-Lipschitz ball coordinates.
The resulting finite-on-compacts maps glue to a locally
bi-Lipschitz map in both directions; use the common polyhedral
coordinates at their interfaces.  Source and target local
finiteness gives a global homeomorphism.  On the cores and tubes
use $h_*$, matching the boundary values.

Finally remake the stairs with their increments allocated only on
$E_j$; Lemma~\ref{wl:stairs} still bounds their root-parameter speeds.
We now pin the point label $q_h(x)$, rather than the support of an
individual routed layer.  Fix a value $q_x=q_h(x)$ and a guard
$\kappa_g$.  Take any point
$z$ on the constant-$q_x$ fibre and any point $y\in\kappa_g$.
The target product-fibre description gives a piecewise linear path
from $h(z)$ to a boundary anchor; pulling it back by $h$ gives
a rectifiable path $\nu_z$ from $z$ to a boundary anchor $a_z$.
For almost every occupied parameter, $\nu_z$ misses its routed
layer: by exact realization, a hit would require $q_x$ to have that
single standard label.  Choose a path $\nu_y\subset\kappa_g$
from $y$ to a boundary anchor $a_g$.  This path misses every
pre-stack, hence every routed layer.  For any rectifiable connector
$\sigma$ from $z$ to $y$, form the boundary-to-boundary path
$\gamma=\nu_z^{-1}*\sigma*\nu_y$.

Use the point-distance potentials for the star and perimeter
coordinates of $q_x$.  The comparison
\eqref{wl:graph-physical} shows that, if the guard label is separated
from $q_x$, one of these potentials separates the two labels.
Its value at $q_x$ is zero before applying the stairs.  Sufficiently
small stair displacement keeps its value at the first transformed
anchor small and its value at the guard anchor bounded away from
zero.  Apply the period identity to $\gamma$ and weight its crossings by
the occupied-stair coefficients as in
Lemma~\ref{wl:crossing-bound}.  The two appended fibre paths contribute
zero for almost every parameter with nonzero coefficient.  Thus the
potential difference at $a_z,a_g$ is bounded by the weighted unsigned
crossings on $\sigma$, and hence by $C\int_\sigma G$.
The coefficient bound changes only by $1+O(\zeta)$.
The same invariant-support tests for remote subedges and vertices
make these prescriptions locally consistent.

If $q_h(x)$ differs too far from $q_*(x)$, follow a path in its
constant-$q_x$ fibre toward an exact boundary anchor.  Since $q_h$
is constant along this path and
$q_*$ agrees with it at the anchor, continuity of $q_*$ produces a
fixed-diameter stretch with a definite discrepancy.  Pack small
balls along this stretch and use nearby guards.  The connector
bound just proved applies to every spherical path between the fibre
and a guard.  The sphere argument of Proposition~\ref{wl:barrier}
therefore gives \eqref{wl:barrier-divergence}, contradicting the same
previously scheduled capacity.  The initial guard spacings accommodate
both the earlier layer-support tests and these point-label tests.
This proves the requested accuracy of $q_h$; the argument estimates
labels after $T$, not $Dh$.
\end{proof}

\subsection{Target stairs, parameter costs, and strict approximation}

Fix the homeomorphism $h$ and occupied intervals supplied by
Theorem~\ref{wl:realization}.  Thus $h$ sends every occupied root
parameter $\lambda\in E_j$ to its standard label $p_j(\lambda)$,
agrees with the prescribed map on full blocks, and respects the
central tubes.  We now collapse the unused target label intervals.
This makes the derivative of the resulting singular composition
depend on the prescribed scalar parameters.

The activation radii and full-cell layer widths are fixed target
data; the occupied stair intervals are supplied by $h$.  Denote the
radial and perimeter stairs by $S_D$ and $S_e$.  For a fixed activation
radius $a$ with $r_{u,D}\ll a\ll1$, use a nondecreasing Lipschitz
cutoff $H_a$ with an absolute Lipschitz bound,
\[
 H_a(u)=0\ (u\leq a),\qquad H_a(u)=u\ (u\geq3a),\qquad
 0\leq H_a(u)\leq u,
\]
with the choices made locally feature by feature.  On the output
two-complex define
\begin{align}
 P_0\bigl(d_D+r(\xi_e(s)-d_D)\bigr)
   &=d_D+\rho_D(r)\bigl(\xi_e(S_e(s))-d_D\bigr),
       \label{wl:target-map}\\
 \rho_D(r)&=H_a\bigl(\exp(S_D(L_D\log r)/L_D)\bigr)
       \quad(r\geq r_{u,D}),\label{wl:radial-stair}
\end{align}
with $\rho_D=0$ farther inward.  The shared subedge stairs make these
formulas agree across sectors and subcells.  By making the stair
accuracy small relative to $L_D$, one has
\begin{equation}\label{wl:target-speeds}
 \rho_D(r)\leq Cr,\qquad |\rho_D'(r)|\leq C/c_*.
\end{equation}
The derivative of the output with respect to one occupied root
parameter is $O(r)$.  Above the central transition, its output
log-radius or perimeter parameter has speed at most $1+O(\zeta)$.
The first bound cancels the polar $1/r$ appearing in weak coarea
estimates when the true and final labels are pinned.  No ratio of
subedge count to subedge length enters these speeds.

Extend $P_0$ to a full convex output cell using a fixed interior point
$z_i$.  Write a point as $z_i+b(y-z_i)$ with $y$ on the cell boundary.
Let $\lambda(b)$ be zero away from a thin boundary layer and rise to
one at $b=1$.  On the boundary, let $\mathcal P_\lambda$ interpolate
the radial factor between $r$ and $\rho_D(r)$ and the perimeter
factor between $s$ and $S_e(s)$, separately and linearly.  Set
\begin{equation}\label{wl:conical-extension}
 P_0\bigl(z_i+b(y-z_i)\bigr)
   =z_i+b\bigl(\mathcal P_{\lambda(b)}(y)-z_i\bigr).
\end{equation}
Here $P_0$ denotes the final singular map; $\mathcal P_0$ is the
identity boundary map.  This notation avoids identifying those two
uses of zero.  The interpolation preserves each boundary sector and
subcell.

Tangential Lipschitz bounds in \eqref{wl:conical-extension} depend only
on the fixed cell geometry, the subcell aspect bounds and $c_*$.  The
speed in $\lambda$ tends uniformly to zero with the radial cutoff
scale and stair displacement.  In fixing these target data, first choose the boundary-layer
thickness using the fixed tangential bounds and absolute continuity
of the relevant integrals; then make the factor displacements small
compared with that thickness.  Thus the extra full-cell energy of
$P_0Th$ can be made summably small.  This uses that $h=h_*$ on the
full input blocks and that $T$ is the prescribed homothety there.
On full blocks outside these thin boundary layers the composition
remains $q_0$.

\begin{proposition}[Parameter cost after realization]\label{wl:parameter-cost}
On occupied pre-stack pieces, a final active root parameter has
local differential $\pm du_j$.  On a routed circle box it has local
differential
\begin{equation}\label{wl:diagonal-gradient}
 \pm d(w_j-u_j+u_k)=\pm(-du_j+du_k),
\end{equation}
with the local coordinate reversals understood.  At such a switched
point only one root scalar is active.  At an essential box both
original root scalars persist independently.  Consequently the
exterior derivative costs of $P_0Th$ depend on these parameter
differentials and the speeds in \eqref{wl:target-speeds}, with no
factor involving $Dh$ or the number of switches.
\end{proposition}

\begin{proof}
Every partial itinerary is a composition of translations and
reflections of the root parameter.  Thus its derivative is $+1$
or $-1$, irrespective of its length.  A diagonal segment has one
reachable ancestry, so there is only one active scalar on the
corresponding circle annulus.  Essential boxes were never switched
and retain both root parameters.  Under the realization, each
occupied interval corresponds by the original affine map $p_j$ to
its target label.  The subsequent stair has bounded forward speed
against that root parameter by Lemma~\ref{wl:stairs}.  In inactive
label directions the stair is constant.  The chain rule in the
piecewise product coordinates therefore yields the asserted local
formulas and the $O(r)$ output factors; interfaces have measure
zero.  This accounts for the derivative entirely in terms of the
original parameter functions.  The qualitative disk and ball
extension supplies no additional energy constant.
\end{proof}

\begin{proposition}[Strict target compositions]\label{wl:strict}
For a fixed realization $h$ as above, the singular composition $P_0Th$
admits locally bi-Lipschitz homeomorphic approximations
$P^{\varepsilon}T_\eta h$ with arbitrarily small prescribed local
$W^{1,p}$ differences and value errors.  Prescriptions may be
summable on a compact exhaustion.  Only local finiteness of the
uncontrolled strict derivatives is required in regions whose energy
is subsequently removed by a source collapse.
\end{proposition}

\begin{proof}
Replace the radial and perimeter factors by
\[
 (1-\varepsilon_D)\rho_D(r)+\varepsilon_Dr,
 \qquad (1-\varepsilon_e)S_e(s)+\varepsilon_es,
\]
with positive, sufficiently small constants on each feature.  Each
factor is strictly increasing with a positive lower slope on a
compact feature.  The interpolating boundary maps in
\eqref{wl:conical-extension} are then homeomorphisms at every radius.
Polar coordinates, including their centers, and the conical cell
coordinates give both local Lipschitz directions.  Hence
$P^{\varepsilon}$ is a locally bi-Lipschitz self-homeomorphism of
$\Lambda$.  Piecewise Lipschitz derivative bounds for the differences
of the factors imply
\[
 P^{\varepsilon}Th\longrightarrow P_0Th
       \quad\hbox{in }W^{1,p}_{\rm loc}
\]
as the feature parameters decrease, with arbitrary compact-local
error prescriptions.

Now fix $P^{\varepsilon}$ and $h$, and use the strict quantizer
$T_\eta$ from Lemma~\ref{qt:strict}.  The composition convergence
requires checking the tangential derivative at a collapsed stratum,
not merely invoking uniform convergence.  Use stairs and cutoffs
that are piecewise smooth, with finitely many pieces on each compact.
On the relative interior of a true face or edge, the pushforward of
source volume by $Th$ is absolutely continuous with respect to the
corresponding dimensional measure: this follows from the product
fibres of $T$ and local bi-Lipschitzness of $h$.  Sector cuts and stair
knots therefore have null inverse image within these strata.

At a generic face point, the derivatives of the full-cell conical
extension on face-tangent vectors converge to the boundary
transformation's derivatives.  At a generic edge point approached
through an incident face sector, an edge-tangent vector keeps $r$
fixed and the perimeter derivative tends to the common subedge
derivative at $r=1$.  In any incident full-cell cone the conical
radial coordinate has derivative zero on that edge-tangent vector;
thus the same compatibility holds on approach through the full
cell.  On a true-vertex level set $D(Th)=0$ almost everywhere.
Consequently the tangential limit supplied by $DT_\eta$ gives
\[
 D(P^{\varepsilon}T_\eta h)
      \longrightarrow D(P^{\varepsilon}Th)
       \quad\hbox{almost everywhere locally}.
\]
The fixed piecewise bounds give an integrable local majorant after
$\eta$ and its gradient are restricted as in
Lemma~\ref{qt:strict}.  Bounded convergence proves the local
power estimates.  Properness and local finiteness allow a diagonal
choice meeting summable compact prescriptions.  Before passage to
the limit, $T_\eta h$ is itself locally bi-Lipschitz, so the ordinary
piecewise chain rule applies.
\end{proof}
